\hypertarget{ux440ux430ux437ux43cux44bux448ux43bux438ux437ux43cux44b-ux43e-ux441ux438ux43dux445ux440ux43eux43dux438ux437ux430ux446ux438ux438-ux444ux430ux437-ux43fux43e-ux43eux431ux440ux430ux437ux430ux43c-ux442ux43eux447ux435ux43a-ux441-commander-sol}{%
\section{Размышлизмы о синхронизации фаз по образам точек с Commander
Sol}\label{ux440ux430ux437ux43cux44bux448ux43bux438ux437ux43cux44b-ux43e-ux441ux438ux43dux445ux440ux43eux43dux438ux437ux430ux446ux438ux438-ux444ux430ux437-ux43fux43e-ux43eux431ux440ux430ux437ux430ux43c-ux442ux43eux447ux435ux43a-ux441-commander-sol}}

\hypertarget{ux442ux43eux447ux43dux44bux435-ux441ux442ux43eux438ux43cux43eux441ux442ux438-ux443ux43fux430ux43aux43eux432ux43aux438-ux43fux440ux43eux441ux442ux440ux430ux43dux441ux442ux432-ux440ux430ux437ux440ux435ux437ux43eux432-ux438-ux442ux43eux447ux43dux44bux439-ux43fux43eux440ux43eux433-ux441ux442ux430ux431ux438ux43bux438ux437ux430ux446ux438ux438-ux432-fcoa}{%
\subsection{Точные стоимости, упаковки пространств разрезов и точный
порог стабилизации в
FCOA}\label{ux442ux43eux447ux43dux44bux435-ux441ux442ux43eux438ux43cux43eux441ux442ux438-ux443ux43fux430ux43aux43eux432ux43aux438-ux43fux440ux43eux441ux442ux440ux430ux43dux441ux442ux432-ux440ux430ux437ux440ux435ux437ux43eux432-ux438-ux442ux43eux447ux43dux44bux439-ux43fux43eux440ux43eux433-ux441ux442ux430ux431ux438ux43bux438ux437ux430ux446ux438ux438-ux432-fcoa}}

\textbf{Alex Malachevsky · Commander Sol}\\
\textbf{Исследовательская рукопись · Версия 1.0-rc1 · 31 августа 2026}

\begin{center}\rule{0.5\linewidth}{0.5pt}\end{center}

\hypertarget{ux430ux43dux43dux43eux442ux430ux446ux438ux44f}{%
\subsection{Аннотация}\label{ux430ux43dux43dux43eux442ux430ux446ux438ux44f}}

Работа выполнена в рамках Fixed-Carrier Oriented Algebra (FCOA),
зафиксированных в \textbf{FCOA Definition 1.0}, \emph{Fixed-Carrier
Oriented Algebra (FCOA): Definition, Typed Partial Operations, Carrier
Erasure, and the Canonical M0 Baseline},
https://doi.org/10.5281/zenodo.22164246. Мы исследуем абстрактную задачу
синхронизации фаз полного носителя, которая возникает после того, как
разреженные компоненты анонимных операций уже допускают корректно
определённые локальные перестановочные фазы. Пусть заданы \texttt{r}
перестановок \texttt{pi\_1,...,pi\_r\ in\ S\_q}. Примитивное ограничение
по образу точки имеет вид \texttt{pi\_i(a)=pi\_j(a)} для одного
исходного цвета \texttt{a}. Через \texttt{L\_q(r)} обозначается
минимальное число таких ограничений, которое заставляет всякий
удовлетворяющий набор перестановок быть диагональным.

Мы доказываем точную графовую переформулировку: система ограничений
синхронизирует тогда и только тогда, когда канонический трансверсальный
quotient-граф имеет единственную правильную \texttt{q}-раскраску с
точностью до глобального переименования цветов. Отсюда следует
необходимая связность объединения любых двух цветовых графов
ограничений. Далее получены бесконечные точные семейства

\[
L_2(r)=r-1,
\qquad
L_3(r)=\left\lceil\frac{3(r-1)}2\right\rceil,
\qquad
L_q(2)=q-1,
\qquad
L_q(3)=2q-3,
\]

и полностью закрыт четырёхфазный столбец:

\[
L_q(4)=
\begin{cases}
3,&q=2,\\
2q-1,&3\le q\le5,\\
12,&q=6,\\
3q-7,&q\ge7.
\end{cases}
\]

Главный структурный результат сопоставляет каждому исходному цвету
нормированное бинарное пространство разрезов
\texttt{W(P\_a)\ \textless{}=\ F\_2\^{}\{r-1\}}, размерность которого
равна экономии числа ограничений для этого цвета. Синхронизация
заставляет ненулевые части этих пространств быть попарно
непересекающимися, поэтому

\[
\sum_a(2^{d_a}-1)\le2^{r-1}-1.
\]

Совпадающая по порядку бинарная конструкция разрезов даёт точный закон
больших алфавитов

\[
\boxed{L_q(r)=(r-1)q-(2^{r-1}-1)}
\]

для всех

\[
q\ge2^{r-1}-1,
\]

причём этот порог доказан точным. Таким образом, при фиксированном
\texttt{r} нерешённая часть задачи конечна. Конечногеометрические
ингредиенты сами по себе классические; заявляемый вклад состоит в модели
FCOA-синхронизации по образам точек, её сведении к unique colorability и
пространствам разрезов, точных семействах и теореме о точной
стабилизации.

\textbf{Ключевые слова:} FCOA; анонимные значения; синхронизация
перестановок; единственная раскраска; разбиения множеств; пространства
разрезов; partial spreads; стоимость жёсткости.

\begin{center}\rule{0.5\linewidth}{0.5pt}\end{center}

\hypertarget{ux432ux432ux435ux434ux435ux43dux438ux435}{%
\subsection{1.
Введение}\label{ux432ux432ux435ux434ux435ux43dux438ux435}}

Анонимные выходные слои FCOA естественно разделяют два вопроса. Первый:
индуцирует ли автоморфизм разреженного реляционного редукта локальную
перестановку видимых имён выходов? Второй начинается лишь после
существования таких локальных фаз: сколько дополнительной информации о
равенствах нужно, чтобы все локальные фазы совпали с одной глобальной
анонимной перестановкой?

Настоящая работа изолирует второй вопрос как конечную экстремальную
задачу.

Пусть

\[
O=\{0,1,\ldots,q-1\}
\]

--- анонимный алфавит, а

\[
\pi_1,\ldots,\pi_r\in S_O.
\]

Для \texttt{1\textless{}=i\textless{}j\textless{}=r} и \texttt{a\ in\ O}
примитивное \textbf{ограничение по образу точки} имеет вид

\[
[i,j;a]:\qquad \pi_i(a)=\pi_j(a).
\tag{1}
\]

Множество \texttt{S} таких ограничений называется
\textbf{синхронизирующим}, если всякий удовлетворяющий набор диагонален:

\[
\pi_1=\cdots=\pi_r.
\tag{2}
\]

Определим

\[
\boxed{
L_q(r)=\min\{|S|:S\text{ синхронизирует}\}.
}
\tag{3}
\]

Наивная spanning-tree конструкция сразу даёт

\[
L_q(r)\le(q-1)(r-1),
\tag{4}
\]

поскольку совпадение двух перестановок на \texttt{q-1} исходных точках
заставляет их совпасть полностью. При \texttt{q=2} это точно. При
\texttt{q\textgreater{}=3} биективность позволяет ограничениям на разных
парах компонент взаимодействовать, поэтому дерево далеко не всегда
оптимально.

Главные результаты статьи:

\begin{enumerate}
\def\labelenumi{\arabic{enumi}.}
\tightlist
\item
  точное сведение синхронизации к unique colorability специального
  трансверсального quotient-графа;
\item
  универсальное условие связности объединения любых двух цветовых графов
  и нижняя граница \[
  L_q(r)\ge\left\lceil\frac{q(r-1)}2\right\rceil;
  \tag{5}
  \]
\item
  точные формулы для всей строки \texttt{q=3} и столбцов
  \texttt{r=2,3,4};
\item
  универсальная бинарная cut-конструкция по всем ненулевым векторам
  \texttt{F\_2\^{}\{r-1\}};
\item
  packing-неравенство для пространств разрезов, совпадающее с
  конструкцией и дающее точную стабилизацию \[
  L_q(r)=(r-1)q-(2^{r-1}-1)
  \quad(q\ge2^{r-1}-1),
  \tag{6}
  \] с точным порогом.
\end{enumerate}

Конечная предстабилизационная область при \texttt{r\textgreater{}=5}
остаётся открытой и отделяется от доказанного хвоста.

\begin{center}\rule{0.5\linewidth}{0.5pt}\end{center}

\hypertarget{ux440ux430ux43cux43aux430-fcoa-ux438-ux43eux431ux43bux430ux441ux442ux44c-ux43fux440ux438ux43cux435ux43dux438ux43cux43eux441ux442ux438}{%
\subsection{2. Рамка FCOA и область
применимости}\label{ux440ux430ux43cux43aux430-fcoa-ux438-ux43eux431ux43bux430ux441ux442ux44c-ux43fux440ux438ux43cux435ux43dux438ux43cux43eux441ux442ux438}}

\hypertarget{ux444ux43eux440ux43cux430ux43bux44cux43dux430ux44f-ux440ux430ux43cux43aux430}{%
\subsubsection{2.1 Формальная
рамка}\label{ux444ux43eux440ux43cux430ux43bux44cux43dux430ux44f-ux440ux430ux43cux43aux430}}

Каноническая рамка задаётся FCOA Definition 1.0 {[}1{]}. Мотивация
исходит из конечного carrier \texttt{G} с частичной операцией,
терминальные выходы которой образуют анонимные value-fibers. Разреженные
реляционные редукты могут разбивать геометрию определённых клеток на
comparison components. В \textbf{full-support phase-admissible sector}
каждая релевантная компонента несёт локальную перестановку одного и того
же анонимного терминального алфавита \texttt{O}.

Настоящая статья переходит к производному слою синхронизации. Носитель
индексов фаз:

\[
R=\{1,\ldots,r\},
\]

терминальный сорт --- анонимное множество \texttt{O}. Примитивные данные
здесь --- не новые значения операции и не новые реальные FCOA-клетки, а
абстрактные равенства (1) между образами одного исходного цвета в двух
локальных фазах. Относительно предыдущего слоя sparse multicolor
transport добавляется только выбранное семейство таких point-image
равенств.

\hypertarget{erasure-ux438-recovery}{%
\subsubsection{2.2 Erasure и recovery}\label{erasure-ux438-recovery}}

\textbf{Erasure convention:} имена терминальных значений остаются
анонимными. Одновременная левая композиция

\[
\pi_i\mapsto\sigma\circ\pi_i
\qquad(\sigma\in S_q)
\tag{7}
\]

меняет только общее именование выходов и сохраняет все ограничения.
Поэтому одну фазу можно нормировать к тождественной.

\textbf{Recovery target:} восстанавливаются не именованные цвета, а одна
глобальная анонимная перестановка, то есть диагональность (2).

\hypertarget{firewall}{%
\subsubsection{2.3 Firewall}\label{firewall}}

Величина \texttt{L\_q(r)} --- абстрактная стоимость синхронизации
full-support фаз. Это не минимум реальных operation cells, требуемых для
ремонта FCOA carrier. Здесь не определяется multicolor-аналог
actual-cell параметра \texttt{alpha\_q} и не утверждаются неравенства,
связывающие такой параметр с \texttt{L\_q(r)}. Арифметика на носителе
FCOA не импортируется.

\begin{center}\rule{0.5\linewidth}{0.5pt}\end{center}

\hypertarget{ux446ux432ux435ux442ux43eux432ux44bux435-ux433ux440ux430ux444ux44b-ux43eux433ux440ux430ux43dux438ux447ux435ux43dux438ux439-ux438-forest-reduction}{%
\subsection{3. Цветовые графы ограничений и forest
reduction}\label{ux446ux432ux435ux442ux43eux432ux44bux435-ux433ux440ux430ux444ux44b-ux43eux433ux440ux430ux43dux438ux447ux435ux43dux438ux439-ux438-forest-reduction}}

Для каждого исходного цвета \texttt{a\ in\ O} определим граф

\[
\Gamma_a=\Gamma_a(S)
\]

на множестве вершин \texttt{R} условием

\[
\{i,j\}\in E(\Gamma_a)
\quad\Longleftrightarrow\quad
[i,j;a]\in S.
\tag{8}
\]

Равенство распространяется вдоль путей, поэтому существенна лишь
компонентная структура \texttt{Gamma\_a}.

\hypertarget{ux43bux435ux43cux43cux430-3.1-forest-reduction}{%
\subsubsection{Лемма 3.1 --- forest
reduction}\label{ux43bux435ux43cux43cux430-3.1-forest-reduction}}

Любую систему ограничений можно заменить системой с тем же множеством
решений и не большим числом ограничений, в которой каждый
\texttt{Gamma\_a} является лесом.

\textbf{Доказательство.} В каждой компоненте связности \texttt{Gamma\_a}
заменим рёбра произвольным остовным деревом. Индуцируемое отношение
равенства значений \texttt{pi\_i(a)} не изменится. Повторяем независимо
для всех \texttt{a}. \texttt{square}

Для редуцированной системы положим

\[
m_a=|E(\Gamma_a)|,
\qquad
c_a=\kappa(\Gamma_a)=r-m_a.
\tag{9}
\]

Тогда

\[
|S|=\sum_a m_a.
\tag{10}
\]

\begin{center}\rule{0.5\linewidth}{0.5pt}\end{center}

\hypertarget{ux441ux438ux43dux445ux440ux43eux43dux438ux437ux430ux446ux438ux44f-ux43aux430ux43a-ux437ux430ux434ux430ux447ux430-ux435ux434ux438ux43dux441ux442ux432ux435ux43dux43dux43eux439-ux440ux430ux441ux43aux440ux430ux441ux43aux438}{%
\subsection{4. Синхронизация как задача единственной
раскраски}\label{ux441ux438ux43dux445ux440ux43eux43dux438ux437ux430ux446ux438ux44f-ux43aux430ux43a-ux437ux430ux434ux430ux447ux430-ux435ux434ux438ux43dux441ux442ux432ux435ux43dux43dux43eux439-ux440ux430ux441ux43aux440ux430ux441ux43aux438}}

Пусть \texttt{B\_a} --- множество компонент связности \texttt{Gamma\_a}.
Построим граф \texttt{H(S)}.

Его вершины имеют вид

\[
(a,B),\qquad a\in O,\ B\in B_a.
\tag{11}
\]

Для каждого индекса фазы \texttt{i} обозначим через \texttt{B\_a(i)}
компоненту \texttt{Gamma\_a}, содержащую \texttt{i}. Соединим попарно
все \texttt{q} вершин

\[
(a,B_a(i)),\qquad a\in O.
\tag{12}
\]

Каждая фаза тем самым задаёт канонический \texttt{K\_q}-трансверсал.
Каноническая правильная раскраска:

\[
\chi_0(a,B)=a.
\tag{13}
\]

\hypertarget{ux442ux435ux43eux440ux435ux43cux430-4.1-ux44dux43aux432ux438ux432ux430ux43bux435ux43dux442ux43dux43eux441ux442ux44c-ux441ux438ux43dux445ux440ux43eux43dux438ux437ux430ux446ux438ux438-ux438-unique-colorability}{%
\subsubsection{Теорема 4.1 --- эквивалентность синхронизации и unique
colorability}\label{ux442ux435ux43eux440ux435ux43cux430-4.1-ux44dux43aux432ux438ux432ux430ux43bux435ux43dux442ux43dux43eux441ux442ux44c-ux441ux438ux43dux445ux440ux43eux43dux438ux437ux430ux446ux438ux438-ux438-unique-colorability}}

Семейство ограничений \texttt{S} синхронизирует тогда и только тогда,
когда \texttt{chi\_0} --- единственная правильная \texttt{q}-раскраска
\texttt{H(S)} с точностью до глобальной перестановки названий цветов.

\textbf{Доказательство.} Пусть \texttt{(pi\_1,...,pi\_r)} удовлетворяет
\texttt{S}. Если \texttt{i,j} лежат в одной компоненте \texttt{B} графа
\texttt{Gamma\_a}, то по путям \texttt{pi\_i(a)=pi\_j(a)}. Поэтому

\[
F(a,B):=\pi_i(a)\qquad(i\in B)
\tag{14}
\]

корректно определено. На трансверсале (12), соответствующем \texttt{i},
значения \texttt{F} равны

\[
\pi_i(0),\ldots,\pi_i(q-1),
\]

и попарно различны. Значит, \texttt{F} --- правильная
\texttt{q}-раскраска.

Обратно, пусть \texttt{F} --- правильная \texttt{q}-раскраска. Каждый
канонический трансверсал есть \texttt{K\_q}, поэтому его вершины
получают все \texttt{q} цветов ровно по одному разу. Определим

\[
\pi_i(a)=F(a,B_a(i)).
\tag{15}
\]

Тогда каждая \texttt{pi\_i} --- перестановка. Если
\texttt{{[}i,j;a{]}\ in\ S}, вершины \texttt{i,j} лежат в одной
компоненте \texttt{Gamma\_a}, и (15) даёт \texttt{pi\_i(a)=pi\_j(a)}.
Таким образом, правильные раскраски находятся во взаимно однозначном
соответствии с удовлетворяющими наборами фаз.

Диагональные наборы --- в точности раскраски, получаемые из
\texttt{chi\_0} одной глобальной перестановкой выходных цветов.
Следовательно, синхронизация эквивалентна единственности \texttt{chi\_0}
с точностью до глобального relabeling. \texttt{square}

Это связывает задачу с классической теорией uniquely colorable graphs,
но quotient-граф сильно ограничен: он строится из \texttt{r}
канонических \texttt{K\_q}-трансверсалов, а идентификации разрешены
только внутри одного исходного цветового класса.

\begin{center}\rule{0.5\linewidth}{0.5pt}\end{center}

\hypertarget{ux441ux432ux44fux437ux43dux43eux441ux442ux44c-ux43fux43eux43fux430ux440ux43dux44bux445-ux43eux431ux44aux435ux434ux438ux43dux435ux43dux438ux439}{%
\subsection{5. Связность попарных
объединений}\label{ux441ux432ux44fux437ux43dux43eux441ux442ux44c-ux43fux43eux43fux430ux440ux43dux44bux445-ux43eux431ux44aux435ux434ux438ux43dux435ux43dux438ux439}}

\hypertarget{ux442ux435ux43eux440ux435ux43cux430-5.1}{%
\subsubsection{Теорема
5.1}\label{ux442ux435ux43eux440ux435ux43cux430-5.1}}

Если \texttt{S} синхронизирует, то для любых двух различных исходных
цветов \texttt{a,b}

\[
\boxed{\Gamma_a\cup\Gamma_b\text{ связен}.}
\tag{16}
\]

\textbf{Доказательство.} Пусть \texttt{Gamma\_a\ union\ Gamma\_b}
несвязен. Выберем непустое собственное объединение его компонент
\texttt{X}. Через разрез \texttt{(X,R\textbackslash{}X)} не проходит ни
одного ограничения цветов \texttt{a,b}. Положим \texttt{sigma=(a\ b)} и

\[
\pi_i=\sigma\quad(i\in X),
\qquad
\pi_i=id\quad(i\notin X).
\tag{17}
\]

Внутренние ограничения сравнивают одинаковые перестановки. Каждое
пересекающее разрез ограничение использует цвет, отличный от
\texttt{a,b}, и потому фиксируется \texttt{sigma}. Все ограничения
выполнены, но набор не диагонален. Противоречие. \texttt{square}

Для редуцированной системы отсюда

\[
m_a+m_b\ge r-1.
\tag{18}
\]

Суммирование по всем неупорядоченным парам даёт

\[
(q-1)\sum_a m_a\ge\binom q2(r-1).
\]

\hypertarget{ux441ux43bux435ux434ux441ux442ux432ux438ux435-5.2-half-density-bound}{%
\subsubsection{Следствие 5.2 --- half-density
bound}\label{ux441ux43bux435ux434ux441ux442ux432ux438ux435-5.2-half-density-bound}}

\[
\boxed{
L_q(r)\ge\left\lceil\frac{q(r-1)}2\right\rceil.
}
\tag{19}
\]

Условие связности является специализацией классического свойства
uniquely colorable graphs {[}3{]}. Здесь оно служит первым движком
нижней оценки в ограниченном трансверсальном quotient-классе.

\begin{center}\rule{0.5\linewidth}{0.5pt}\end{center}

\hypertarget{ux442ux43eux447ux43dux44bux435-ux43dux438ux437ux43aux43eux440ux430ux437ux43cux435ux440ux43dux44bux435-ux441ux435ux43cux435ux439ux441ux442ux432ux430}{%
\subsection{6. Точные низкоразмерные
семейства}\label{ux442ux43eux447ux43dux44bux435-ux43dux438ux437ux43aux43eux440ux430ux437ux43cux435ux440ux43dux44bux435-ux441ux435ux43cux435ux439ux441ux442ux432ux430}}

\hypertarget{ux434ux432ux43eux438ux447ux43dux44bux439-ux430ux43bux444ux430ux432ux438ux442}{%
\subsubsection{6.1 Двоичный
алфавит}\label{ux434ux432ux43eux438ux447ux43dux44bux439-ux430ux43bux444ux430ux432ux438ux442}}

При \texttt{q=2} совпадение на одном исходном цвете идентифицирует две
перестановки двухэлементного множества. Поэтому связный граф на
\texttt{r} фазовых индексах необходим и достаточен.

\hypertarget{ux442ux435ux43eux440ux435ux43cux430-6.1}{%
\subsubsection{Теорема
6.1}\label{ux442ux435ux43eux440ux435ux43cux430-6.1}}

\[
\boxed{L_2(r)=r-1.}
\tag{20}
\]

\hypertarget{ux434ux432ux435-ux444ux430ux437ux44b}{%
\subsubsection{6.2 Две
фазы}\label{ux434ux432ux435-ux444ux430ux437ux44b}}

Две перестановки \texttt{q}-элементного множества, совпадающие на
\texttt{q-1} исходных точках, совпадают на последней. Если ограничений
меньше \texttt{q-1}, остаются как минимум две свободные исходные точки,
которые можно транспонировать.

\hypertarget{ux442ux435ux43eux440ux435ux43cux430-6.2}{%
\subsubsection{Теорема
6.2}\label{ux442ux435ux43eux440ux435ux43cux430-6.2}}

\[
\boxed{L_q(2)=q-1.}
\tag{21}
\]

\hypertarget{ux442ux440ux438-ux446ux432ux435ux442ux430}{%
\subsubsection{6.3 Три
цвета}\label{ux442ux440ux438-ux446ux432ux435ux442ux430}}

При \texttt{q=3} из (18)

\[
m_0+m_1\ge r-1,
\quad
m_0+m_2\ge r-1,
\quad
m_1+m_2\ge r-1,
\]

поэтому

\[
L_3(r)\ge\left\lceil\frac{3(r-1)}2\right\rceil.
\tag{22}
\]

Совпадающая конструкция использует трёхconstraintный gadget. Пусть
\texttt{rho} уже синхронизирована, а \texttt{sigma,tau\ in\ S\_3} ---
новые фазы. Наложим

\[
[\rho,\sigma;0],
\qquad
[\rho,\tau;1],
\qquad
[\sigma,\tau;2].
\tag{23}
\]

После глобальной нормировки \texttt{rho=id}. Тогда \texttt{sigma}
фиксирует \texttt{0}, \texttt{tau} фиксирует \texttt{1}, и
\texttt{sigma(2)=tau(2)}. Если бы \texttt{sigma} была нетривиальной
перестановкой, фиксирующей \texttt{0}, то \texttt{sigma(2)=1}, что
невозможно для \texttt{tau(2)}, поскольку \texttt{tau(1)=1}. Значит,
\texttt{sigma=id}; тогда \texttt{tau(2)=2}, откуда \texttt{tau=id}.

Добавляя фазы попарно, получаем точную границу.

\hypertarget{ux442ux435ux43eux440ux435ux43cux430-6.3}{%
\subsubsection{Теорема
6.3}\label{ux442ux435ux43eux440ux435ux43cux430-6.3}}

Для всех \texttt{r\textgreater{}=1}

\[
\boxed{
L_3(r)=\left\lceil\frac{3(r-1)}2\right\rceil.
}
\tag{24}
\]

\hypertarget{ux442ux440ux438-ux444ux430ux437ux44b}{%
\subsubsection{6.4 Три
фазы}\label{ux442ux440ux438-ux444ux430ux437ux44b}}

При \texttt{r=3} каждый редуцированный \texttt{Gamma\_a} имеет
\texttt{m\_a\ in\ \{0,1,2\}}. Если один цвет имеет \texttt{m\_a=0},
всякий другой цвет должен быть связным, и стоимость не меньше
\texttt{2(q-1)}. Иначе каждый цвет имеет как минимум одно ребро. Два
одно-рёберных цвета должны использовать разные рёбра треугольника фаз,
иначе их объединение несвязно. Рёбер всего три, поэтому не более трёх
цветов имеют \texttt{m\_a=1}, а остальные --- \texttt{m\_a=2}.
Следовательно,

\[
L_q(3)\ge3+2(q-3)=2q-3.
\tag{25}
\]

Для верхней оценки используем трёхцветный triangle gadget

\[
[0,1;0],\quad[0,2;1],\quad[1,2;2]
\tag{26}
\]

и для каждого дополнительного цвета \texttt{a\textgreater{}=3}
ограничения \texttt{{[}0,1;a{]}}, \texttt{{[}0,2;a{]}}. Дополнительные
цвета фиксируются на всех фазах, оставляя действие \texttt{S\_3},
которое убивается (26).

\hypertarget{ux442ux435ux43eux440ux435ux43cux430-6.4}{%
\subsubsection{Теорема
6.4}\label{ux442ux435ux43eux440ux435ux43cux430-6.4}}

Для всех \texttt{q\textgreater{}=3}

\[
\boxed{L_q(3)=2q-3.}
\tag{27}
\]

\begin{center}\rule{0.5\linewidth}{0.5pt}\end{center}

\hypertarget{ux43fux43eux43bux43dux44bux439-ux447ux435ux442ux44bux440ux451ux445ux444ux430ux437ux43dux44bux439-ux441ux442ux43eux43bux431ux435ux446}{%
\subsection{7. Полный четырёхфазный
столбец}\label{ux43fux43eux43bux43dux44bux439-ux447ux435ux442ux44bux440ux451ux445ux444ux430ux437ux43dux44bux439-ux441ux442ux43eux43bux431ux435ux446}}

При \texttt{r=4} каждый исходный цвет задаёт разбиение \texttt{P\_a}
четырёх фаз. После forest reduction его ранг

\[
m_a=4-|P_a|
\]

лежит в \texttt{\{0,1,2,3\}}.

Условие совместимости из Теоремы 5.1 принимает вид

\[
P_a\vee P_b=\mathbf 1
\qquad(a\ne b),
\tag{28}
\]

где \texttt{mathbf\ 1} --- одно-блочное разбиение.

Решётка разбиений четырёх точек даёт полный список совместимости малых
рангов:

\begin{itemize}
\tightlist
\item
  дискретное разбиение ранга ноль совместимо только с одно-блочным
  разбиением ранга три;
\item
  два rank-1 разбиения никогда не совместимы;
\item
  фиксированное rank-1 разбиение совместимо ровно с четырьмя rank-2
  разбиениями;
\item
  rank-2 разбиений ровно семь, и любые два различных совместимы;
\item
  rank-3 разбиение совместимо со всеми.
\end{itemize}

Пусть \texttt{n\_k} --- число цветов ранга \texttt{k}. Если rank-0 цвета
нет,

\[
|S|=n_1+2n_2+3n_3=3q-(2n_1+n_2).
\tag{29}
\]

Классификация даёт максимальный defect

\[
2n_1+n_2=\begin{cases}
q+1,&3\le q\le5,\\
6,&q=6,\\
7,&q\ge7.
\end{cases}
\tag{30}
\]

и соответствующие нижние границы.

При \texttt{q=3} пятиconstraintная конструкция является экземпляром уже
доказанного рекурсивного \texttt{S\_3} gadget Теоремы 6.3 после
переименования фаз и цветов. При \texttt{q=4,5} можно использовать один
rank-1 partition и соответственно три или четыре совместимых rank-2
партнёра; прямая нормировка и инъективность вынуждают все фазы быть
тождественными. При \texttt{q=6} оптимальный \texttt{q=5} gadget
дополняется одним связным цветом. При \texttt{q=7} используются все семь
bipartitions четырёх фаз --- это случай \texttt{r=4} универсальной
binary-cut конструкции следующего раздела. При \texttt{q\textgreater{}7}
новые цвета делаются связными.

\hypertarget{ux442ux435ux43eux440ux435ux43cux430-7.1-ux442ux43eux447ux43dux430ux44f-ux447ux435ux442ux44bux440ux451ux445ux444ux430ux437ux43dux430ux44f-ux441ux442ux43eux438ux43cux43eux441ux442ux44c}{%
\subsubsection{Теорема 7.1 --- точная четырёхфазная
стоимость}\label{ux442ux435ux43eux440ux435ux43cux430-7.1-ux442ux43eux447ux43dux430ux44f-ux447ux435ux442ux44bux440ux451ux445ux444ux430ux437ux43dux430ux44f-ux441ux442ux43eux438ux43cux43eux441ux442ux44c}}

\[
\boxed{
L_q(4)=
\begin{cases}
3,&q=2,\\
2q-1,&3\le q\le5,\\
12,&q=6,\\
3q-7,&q\ge7.
\end{cases}
}
\tag{31}
\]

Изолированное значение при \texttt{q=6} имеет структурную природу:
rank-1 архитектура допускает не более четырёх rank-2 цветов, тогда как
all-bipartition архитектура имеет capacity семь.

\begin{center}\rule{0.5\linewidth}{0.5pt}\end{center}

\hypertarget{ux443ux43dux438ux432ux435ux440ux441ux430ux43bux44cux43dux430ux44f-ux431ux438ux43dux430ux440ux43dux430ux44f-ux441ux438ux43dux445ux440ux43eux43dux438ux437ux430ux446ux438ux44f-ux440ux430ux437ux440ux435ux437ux430ux43cux438}{%
\subsection{8. Универсальная бинарная синхронизация
разрезами}\label{ux443ux43dux438ux432ux435ux440ux441ux430ux43bux44cux43dux430ux44f-ux431ux438ux43dux430ux440ux43dux430ux44f-ux441ux438ux43dux445ux440ux43eux43dux438ux437ux430ux446ux438ux44f-ux440ux430ux437ux440ux435ux437ux430ux43cux438}}

Зафиксируем \texttt{r\textgreater{}=2} и положим

\[
n=r-1,
\qquad
V=\mathbb F_2^n\setminus\{0\}.
\tag{32}
\]

Используем один активный исходный цвет для каждого вектора

\[
v=(v_1,\ldots,v_n)\in V.
\]

Число активных цветов

\[
q_0=|V|=2^{r-1}-1.
\tag{33}
\]

Для цвета \texttt{v} разобьём фазы на

\[
B_v^0=\{0\}\cup\{i:v_i=0\},
\qquad
B_v^1=\{i:v_i=1\}.
\tag{34}
\]

Внутри каждого блока выберем остовное дерево. Поскольку блоков два и оба
непусты, цвет \texttt{v} стоит \texttt{r-2} ограничений.

\hypertarget{ux442ux435ux43eux440ux435ux43cux430-8.1-binary-cut-gadget}{%
\subsubsection{Теорема 8.1 --- binary-cut
gadget}\label{ux442ux435ux43eux440ux435ux43cux430-8.1-binary-cut-gadget}}

Семейство (34) синхронизирует \texttt{r} фаз в
\texttt{S\_\{2\^{}\{r-1\}-1\}}.

\textbf{Доказательство.} Нормируем \texttt{pi\_0=id}. Для
\texttt{i\textgreater{}=1} положим

\[
H_i=\{v:v_i=0\},
\qquad
A_i=\{v:v_i=1\}.
\tag{35}
\]

Если \texttt{v\ in\ H\_i}, фазы \texttt{0} и \texttt{i} лежат в одном
связном блоке \texttt{B\_v\^{}0}, поэтому

\[
\pi_i(v)=v.
\tag{36}
\]

Значит, \texttt{pi\_i} фиксирует \texttt{H\_i} поточечно и сохраняет
\texttt{A\_i} как множество.

Для \texttt{j!=i} и \texttt{v\ in\ A\_i\ cap\ A\_j} фазы \texttt{i,j}
лежат в одном блоке \texttt{B\_v\^{}1}, поэтому

\[
\pi_i(v)=\pi_j(v).
\tag{37}
\]

Левая часть лежит в \texttt{A\_i}, правая --- в \texttt{A\_j}, значит
общее значение лежит в \texttt{A\_i\ cap\ A\_j}. Поэтому \texttt{pi\_i}
сохраняет каждое множество \texttt{A\_i\ cap\ A\_j}.

На \texttt{A\_i} признаки принадлежности множествам
\texttt{A\_i\ cap\ A\_j}, \texttt{j!=i}, являются в точности оставшимися
координатами \texttt{(v\_j)\_\{j!=i\}} и различают все векторы
\texttt{A\_i}. Следовательно, \texttt{pi\_i} фиксирует каждую точку
\texttt{A\_i}, а вместе с (36) получаем \texttt{pi\_i=id}. Это верно для
всех \texttt{i}, следовательно, все фазы диагональны до нормировки.
\texttt{square}

Стоимость gadget равна

\[
(r-2)(2^{r-1}-1).
\tag{38}
\]

При \texttt{q\textgreater{}q\_0} каждый новый цвет делается связным на
всех \texttt{r} фазах за \texttt{r-1} ограничений. Поэтому

\[
L_q(r)\le(r-1)q-(2^{r-1}-1)
\quad(q\ge q_0).
\tag{39}
\]

\begin{center}\rule{0.5\linewidth}{0.5pt}\end{center}

\hypertarget{ux43fux440ux43eux441ux442ux440ux430ux43dux441ux442ux432ux430-ux440ux430ux437ux440ux435ux437ux43eux432-ux438-packing-bound}{%
\subsection{9. Пространства разрезов и packing
bound}\label{ux43fux440ux43eux441ux442ux440ux430ux43dux441ux442ux432ux430-ux440ux430ux437ux440ux435ux437ux43eux432-ux438-packing-bound}}

Совпадающая нижняя граница получается из канонического кодирования
компонентных разбиений каждого исходного цвета.

Пусть \texttt{P} --- разбиение множества фаз

\[
R=\{0,1,\ldots,r-1\}.
\]

Нормируем бинарные разрезы, фиксируя бит фазы \texttt{0} равным нулю, и
отождествим их с

\[
\mathbb F_2^{r-1}.
\]

Определим

\[
W(P)=\{x\in\mathbb F_2^{r-1}:x\text{ постоянно на каждом блоке }P\},
\tag{40}
\]

где подразумевается расширение \texttt{x\_0=0}.

\hypertarget{ux43bux435ux43cux43cux430-9.1}{%
\subsubsection{Лемма 9.1}\label{ux43bux435ux43cux43cux430-9.1}}

Если \texttt{P} имеет \texttt{c} блоков, то

\[
\dim W(P)=c-1.
\tag{41}
\]

\textbf{Доказательство.} Блок, содержащий фазу \texttt{0}, вынужден
иметь бит ноль. Каждый из остальных \texttt{c-1} блоков выбирает бит
независимо. Поэтому \texttt{\textbar{}W(P)\textbar{}=2\^{}\{c-1\}}.
\texttt{square}

Для компонентного разбиения \texttt{P\_a} леса \texttt{Gamma\_a}
определим экономию

\[
d_a=(r-1)-m_a.
\tag{42}
\]

Поскольку \texttt{m\_a=r-c\_a},

\[
d_a=c_a-1=\dim W(P_a).
\tag{43}
\]

\hypertarget{ux43bux435ux43cux43cux430-9.2-ux442ux43eux436ux434ux435ux441ux442ux432ux43e-joinintersection}{%
\subsubsection{Лемма 9.2 --- тождество
join/intersection}\label{ux43bux435ux43cux43cux430-9.2-ux442ux43eux436ux434ux435ux441ux442ux432ux43e-joinintersection}}

Для любых разбиений \texttt{P,Q}

\[
\boxed{W(P)\cap W(Q)=W(P\vee Q).}
\tag{44}
\]

\textbf{Доказательство.} Нормированный бинарный разрез принадлежит обоим
пространствам тогда и только тогда, когда он постоянен на каждом блоке
\texttt{P} и каждом блоке \texttt{Q}. Это эквивалентно постоянству на
классах эквивалентности, порождённых совместно этими блоками, то есть на
блоках \texttt{P\ vee\ Q}. \texttt{square}

Из Теоремы 5.1

\[
P_a\vee P_b=\mathbf 1
\quad(a\ne b),
\]

следовательно,

\[
W(P_a)\cap W(P_b)=\{0\}.
\tag{45}
\]

Ненулевые векторы пространств разрезов попарно не пересекаются. В
\texttt{F\_2\^{}\{r-1\}} ровно \texttt{2\^{}\{r-1\}-1} ненулевых
векторов.

\hypertarget{ux442ux435ux43eux440ux435ux43cux430-9.3-defect-packing-inequality}{%
\subsubsection{Теорема 9.3 --- defect packing
inequality}\label{ux442ux435ux43eux440ux435ux43cux430-9.3-defect-packing-inequality}}

Всякая синхронизирующая система удовлетворяет

\[
\boxed{
\sum_a(2^{d_a}-1)\le2^{r-1}-1.
}
\tag{46}
\]

Сам подсчёт (46) после получения cut-space reduction --- стандартная
конечногеометрическая packing-оценка и не заявляется как новое
утверждение о partial spreads или vector-space partitions {[}5,6{]}.

Так как

\[
d\le2^d-1
\qquad(d\ge0),
\tag{47}
\]

получаем

\[
\sum_a d_a\le2^{r-1}-1.
\tag{48}
\]

Но из (42)

\[
|S|=(r-1)q-\sum_a d_a.
\tag{49}
\]

Поэтому:

\hypertarget{ux441ux43bux435ux434ux441ux442ux432ux438ux435-9.4-ux443ux43dux438ux432ux435ux440ux441ux430ux43bux44cux43dux430ux44f-defect-ux433ux440ux430ux43dux438ux446ux430}{%
\subsubsection{Следствие 9.4 --- универсальная
defect-граница}\label{ux441ux43bux435ux434ux441ux442ux432ux438ux435-9.4-ux443ux43dux438ux432ux435ux440ux441ux430ux43bux44cux43dux430ux44f-defect-ux433ux440ux430ux43dux438ux446ux430}}

\[
\boxed{
L_q(r)\ge(r-1)q-(2^{r-1}-1).
}
\tag{50}
\]

\begin{center}\rule{0.5\linewidth}{0.5pt}\end{center}

\hypertarget{ux442ux43eux447ux43dux430ux44f-ux441ux442ux430ux431ux438ux43bux438ux437ux430ux446ux438ux44f-ux43fux440ux438-ux431ux43eux43bux44cux448ux438ux445-ux430ux43bux444ux430ux432ux438ux442ux430ux445}{%
\subsection{10. Точная стабилизация при больших
алфавитах}\label{ux442ux43eux447ux43dux430ux44f-ux441ux442ux430ux431ux438ux43bux438ux437ux430ux446ux438ux44f-ux43fux440ux438-ux431ux43eux43bux44cux448ux438ux445-ux430ux43bux444ux430ux432ux438ux442ux430ux445}}

Объединяем разделы 8 и 9.

\hypertarget{ux442ux435ux43eux440ux435ux43cux430-10.1-ux442ux43eux447ux43dux430ux44f-ux442ux435ux43eux440ux435ux43cux430-ux441ux442ux430ux431ux438ux43bux438ux437ux430ux446ux438ux438}{%
\subsubsection{Теорема 10.1 --- точная теорема
стабилизации}\label{ux442ux435ux43eux440ux435ux43cux430-10.1-ux442ux43eux447ux43dux430ux44f-ux442ux435ux43eux440ux435ux43cux430-ux441ux442ux430ux431ux438ux43bux438ux437ux430ux446ux438ux438}}

Для всех \texttt{r\textgreater{}=2} и

\[
q\ge2^{r-1}-1
\]

имеем

\[
\boxed{
L_q(r)=(r-1)q-(2^{r-1}-1).
}
\tag{51}
\]

Причём порог точен: если

\[
q<2^{r-1}-1,
\]

то

\[
L_q(r)>(r-1)q-(2^{r-1}-1).
\tag{52}
\]

\textbf{Доказательство.} Нижняя граница --- Следствие 9.4. Binary-cut
construction даёт равенство при \texttt{q\textgreater{}=2\^{}\{r-1\}-1}.

Предположим, что равенство в (50) достигнуто. Тогда

\[
\sum_a d_a=2^{r-1}-1.
\tag{53}
\]

Из (46) и (47)

\[
\sum_a d_a
\le
\sum_a(2^{d_a}-1)
\le
2^{r-1}-1.
\tag{54}
\]

Равенство в (53) вынуждает равенство в каждом неравенстве
\texttt{d\_a\textless{}=2\^{}\{d\_a\}-1}, что для неотрицательных целых
возможно только при \texttt{d\_a\ in\ \{0,1\}}. Значит, каждый цвет с
положительным defect даёт ровно один ненулевой cut-вектор. Чтобы
получить суммарный defect \texttt{2\^{}\{r-1\}-1}, нужно как минимум
столько же цветов. Следовательно,

\[
q\ge2^{r-1}-1.
\]

Ниже порога равенство невозможно. \texttt{square}

\hypertarget{ux441ux43bux435ux434ux441ux442ux432ux438ux435-10.2}{%
\subsubsection{Следствие
10.2}\label{ux441ux43bux435ux434ux441ux442ux432ux438ux435-10.2}}

Точный порог стабилизации:

\[
\boxed{q_0(r)=2^{r-1}-1.}
\tag{55}
\]

Например,

\[
L_q(5)=4q-15\qquad(q\ge15),
\tag{56}
\]

и

\[
L_q(6)=5q-31\qquad(q\ge31).
\tag{57}
\]

При фиксированном числе фаз вне линейного точного хвоста остаётся лишь
конечное число алфавитов.

\begin{center}\rule{0.5\linewidth}{0.5pt}\end{center}

\hypertarget{ux43aux43eux43dux435ux447ux43dux430ux44f-ux43fux440ux435ux434ux441ux442ux430ux431ux438ux43bux438ux437ux430ux446ux438ux43eux43dux43dux430ux44f-ux437ux430ux434ux430ux447ux430}{%
\subsection{11. Конечная предстабилизационная
задача}\label{ux43aux43eux43dux435ux447ux43dux430ux44f-ux43fux440ux435ux434ux441ux442ux430ux431ux438ux43bux438ux437ux430ux446ux438ux43eux43dux43dux430ux44f-ux437ux430ux434ux430ux447ux430}}

Доказательство Теоремы 10.1 даёт более сильный интерфейс для нерешённого
сектора. Определим

\[
\mathcal D_r(q)
=
\max\left\{
\sum_{i=1}^q\dim W(P_i):
W(P_i)\cap W(P_j)=\{0\}\ (i\ne j)
\right\},
\tag{58}
\]

где нульмерные пространства разрешены для дополнения семейства.

Тогда необходимо

\[
L_q(r)\ge(r-1)q-\mathcal D_r(q).
\tag{59}
\]

Это лишь lower-bound interface: попарно тривиальное пересечение
необходимо для синхронизации, но само по себе может быть недостаточно
для unique colorability трансверсального quotient-графа.

Действительно открытый режим:

\[
\boxed{
4\le q<2^{r-1}-1,
\qquad r\ge5.
}
\tag{60}
\]

Для \texttt{r=5} остаются только \texttt{4\textless{}=q\textless{}=14}.
Точные weighted-packing вычисления полезны как hostile check, но их
значения не повышаются здесь до LQR-теорем без отдельного доказательства
синхронизации.

\begin{center}\rule{0.5\linewidth}{0.5pt}\end{center}

\hypertarget{ux43aux43eux43dux435ux447ux43dux430ux44f-ux43fux440ux43eux432ux435ux440ux43aux430}{%
\subsection{12. Конечная
проверка}\label{ux43aux43eux43dux435ux447ux43dux430ux44f-ux43fux440ux43eux432ux435ux440ux43aux430}}

В репозитории поддерживаются два независимых verifier-скрипта.

\texttt{verify\_lqr.py} проверяет:

\begin{itemize}
\tightlist
\item
  точные \texttt{q=3} конструкции для малых \texttt{r} нормированным
  полным перебором;
\item
  точные \texttt{r=3} конструкции для малых \texttt{q};
\item
  первые значения \texttt{r=4} через partition quotient.
\end{itemize}

\texttt{verify\_lqr\_cutspace.py} независимо проверяет:

\begin{itemize}
\tightlist
\item
  Bell-number enumeration разбиений фаз до \texttt{r=6};
\item
  \texttt{\textbar{}W(P)\textbar{}=2\^{}\{\textbar{}P\textbar{}-1\}};
\item
  \texttt{W(P)\ cap\ W(Q)=W(P\ vee\ Q)} для всех пар разбиений до
  \texttt{r=6};
\item
  weighted partition-packing capacity при \texttt{r=5} до порога
  \texttt{q=15}.
\end{itemize}

Доказательства бесконечных формул от этих вычислений не зависят.

\begin{center}\rule{0.5\linewidth}{0.5pt}\end{center}

\hypertarget{ux43bux438ux442ux435ux440ux430ux442ux443ux440ux43dux430ux44f-ux43fux43eux437ux438ux446ux438ux44f-ux438-ux433ux440ux430ux43dux438ux446ux430-ux43dux43eux432ux438ux437ux43dux44b}{%
\subsection{13. Литературная позиция и граница
новизны}\label{ux43bux438ux442ux435ux440ux430ux442ux443ux440ux43dux430ux44f-ux43fux43eux437ux438ux446ux438ux44f-ux438-ux433ux440ux430ux43dux438ux446ux430-ux43dux43eux432ux438ux437ux43dux44b}}

Графовая и конечногеометрическая окружающая теория классическая.

Harary, Hedetniemi и Robinson {[}3{]} установили фундаментальные
свойства uniquely colorable graphs, включая связность подграфа,
индуцированного любой парой цветовых классов. Bollobas {[}4{]} развил
дальнейшие структурные результаты. Поэтому Теорема 5.1 не заявляется как
новое общее утверждение unique-colorability; она используется как
специализированная форма, индуцированная LQR transversal quotient.

Попарно тривиально пересекающиеся подпространства, partial spreads и
vector-space partitions также классические {[}5,6{]}. После (45) подсчёт
их ненулевых векторов стандартен. Поэтому novelty claim сознательно
ограничен следующей цепочкой внутри задачи FCOA phase synchronization:

\[
\text{point-image constraints}
\longrightarrow
\text{transversal unique-coloring quotient}
\longrightarrow
\text{component partitions}
\longrightarrow
\text{canonical binary cut spaces}
\longrightarrow
\text{sharp LQR stabilization}.
\tag{61}
\]

Permutation synchronization имеет развитую алгоритмическую литературу,
обычно ориентированную на восстановление глобально согласованных
matching-перестановок из шумных попарных оценок {[}7{]}. Здесь параметр
иной: минимальное число точных same-source point-image равенств,
достаточных, чтобы любой набор в \texttt{S\_q\^{}r} стал диагональным.

Специализированный поиск в рамках проекта не обнаружил точного
экстремального параметра \texttt{L\_q(r)} именно в этой форме. Это
консервативная формулировка, а не абсолютное заявление приоритета.

\begin{center}\rule{0.5\linewidth}{0.5pt}\end{center}

\hypertarget{ux43eux433ux440ux430ux43dux438ux447ux435ux43dux438ux44f-ux438-ux43eux442ux43aux440ux44bux442ux44bux435-ux432ux43eux43fux440ux43eux441ux44b}{%
\subsection{14. Ограничения и открытые
вопросы}\label{ux43eux433ux440ux430ux43dux438ux447ux435ux43dux438ux44f-ux438-ux43eux442ux43aux440ux44bux442ux44bux435-ux432ux43eux43fux440ux43eux441ux44b}}

Основные ограничения явные.

\begin{enumerate}
\def\labelenumi{\arabic{enumi}.}
\tightlist
\item
  \texttt{L\_q(r)} измеряет абстрактные phase-link constraints, а не
  реальные FCOA operation-cell additions.
\item
  Предстабилизационный сектор (60) в общем случае не решён.
\item
  Pairwise cut-space packing --- необходимое условие и может не
  исчерпывать все unique-coloring obstructions.
\item
  Multicolor-аналог binary actual-cell repair parameter здесь не
  определяется.
\item
  Существование permutation-valued local phases предполагается;
  разреженные ternary reducts при \texttt{q\textgreater{}=3} могут
  разрушаться до стадии синхронизации, что рассматривается в companion
  manuscript {[}2{]}.
\end{enumerate}

Естественные следующие задачи: определить \texttt{L\_4(r)} для общего
\texttt{r}, решить \texttt{L\_q(5)} при
\texttt{4\textless{}=q\textless{}=14} и охарактеризовать, когда
оптимальная partition-subspace packing конфигурация реально задаёт
синхронизирующий transversal quotient.

\begin{center}\rule{0.5\linewidth}{0.5pt}\end{center}

\hypertarget{ux437ux430ux43aux43bux44eux447ux435ux43dux438ux435}{%
\subsection{15.
Заключение}\label{ux437ux430ux43aux43bux44eux447ux435ux43dux438ux435}}

Стоимость синхронизации по образам точек имеет существенно более тонкую
структуру, чем наивная spanning-tree оценка. Правильные нижнеоценочные
объекты --- не просто графы на фазовых индексах, а компонентные
разбиения; после нормировки эти разбиения порождают бинарные cut-spaces,
ненулевые векторы которых обязаны упаковываться без пересечений.

Это даёт три слоя точной структуры:

\[
L_3(r)=\left\lceil\frac{3(r-1)}2\right\rceil,
\]

точный четырёхфазный переход

\[
L_q(4)=3,\ 2q-1,\ 12,\ 3q-7
\]

на соответствующих диапазонах и общую теорему стабилизации

\[
\boxed{
L_q(r)=(r-1)q-(2^{r-1}-1)
\quad\text{при достижимости линейной границы, причём достижимость начинается ровно при }q=2^{r-1}-1.
}
\]

Следовательно, при фиксированном числе фаз нерешённая задача конечна. В
терминах FCOA это точная capacity theorem для синхронизации уже
существующих анонимных full-support фаз, при сохранении строгого
различия с более трудной задачей реализации синхронизации реальными
operation-cell extensions.

\begin{center}\rule{0.5\linewidth}{0.5pt}\end{center}

\hypertarget{ux43bux438ux442ux435ux440ux430ux442ux443ux440ux430}{%
\subsection{Литература}\label{ux43bux438ux442ux435ux440ux430ux442ux443ux440ux430}}

{[}1{]} A. Malachevsky and Commander Sol, \textbf{Fixed-Carrier Oriented
Algebra (FCOA): Definition, Typed Partial Operations, Carrier Erasure,
and the Canonical M0 Baseline}, 2026. DOI:
\texttt{10.5281/zenodo.22164246}.
https://doi.org/10.5281/zenodo.22164246.

{[}2{]} A. Malachevsky and Commander Sol, \textbf{Reflections on Sparse
Multicolor Transport with Commander Sol: Proper-Coloring Obstructions,
Local Phase Groupoids, and Exact Gluing in FCOA}, companion manuscript,
2026, repository package
\texttt{papers/FCOA-SPARSE-MULTICOLOR-TRANSPORT/}.

{[}3{]} F. Harary, S. T. Hedetniemi, R. W. Robinson, \textbf{Uniquely
colorable graphs}, \emph{Journal of Combinatorial Theory} 6(3) (1969),
264--270. DOI: \texttt{10.1016/S0021-9800(69)80086-4}.

{[}4{]} B. Bollobas, \textbf{Uniquely colorable graphs}, \emph{Journal
of Combinatorial Theory, Series B} 25 (1978), 54--61. DOI:
\texttt{10.1016/S0095-8956(78)80010-0}.

{[}5{]} O. Heden, \textbf{A survey of the different types of vector
space partitions}, \emph{Discrete Mathematics, Algorithms and
Applications} 4(1) (2012), 1250001. DOI:
\texttt{10.1142/S1793830912500012}.

{[}6{]} T. Honold, M. Kiermaier, S. Kurz, \textbf{Partial spreads and
vector space partitions}, in \emph{Network Coding and Subspace Designs},
Springer, 2018, 131--170. DOI: \texttt{10.1007/978-3-319-70293-3\_7}.

{[}7{]} D. Pachauri, T. Kondor, V. Singh, \textbf{Solving the multi-way
matching problem by permutation synchronization}, \emph{Advances in
Neural Information Processing Systems 26}, 2013.

\begin{center}\rule{0.5\linewidth}{0.5pt}\end{center}

\hypertarget{ux432ux43eux441ux43fux440ux43eux438ux437ux432ux43eux434ux438ux43cux43eux441ux442ux44c}{%
\subsection{Воспроизводимость}\label{ux432ux43eux441ux43fux440ux43eux438ux437ux432ux43eux434ux438ux43cux43eux441ux442ux44c}}

Исходные теоремные файлы и независимые конечные verifier-скрипты
поддерживаются в

\texttt{delegated/FCOA\_RIGIDITY\_COST/QGE3/}

репозитория

\texttt{https://github.com/AIDevelopersMonster/Riemann-Hypothesis-Commander-Sol}

в исследовательской ветке \texttt{director/fcoa-rigidity-cost} на момент
сборки рукописи.
